\chapter{Foreword}


The work of this thesis is focused on the study of light scattering in random correlated systems, on one side, and in magneto-optical media, on the other.\\
Since the foundations of the electromagnetic theory, efforts have been made to understand how light interacts with complex environments. Until now, most works are centred in either perfectly ordered or disordered systems. In these extreme cases a rich phenomenology was found, such as the existence of photonic bandgaps in photonic crystals or weak and strong localisation in disordered systems. Between these two limits, perfectly ordered and disordered systems, are the random correlated systems. Although it is known that correlations in disordered media can play an important role in the wave transport properties, these random correlated systems remain largely unexplored. Examples such as the conductivity of liquid metals in solid state physics, or light transport properties in photonic systems governed by short-range structural correlations, like in cornea, demonstrate the importance of these media.\\
Being dominated by multiple light scattering, light transport and emission are closely related phenomena. However, much more attention has been given to the first over the second. In recent years a small number of works have shown the importance of light emission by point sources for the study of correlations in random systems. In particular, these take advantage of the proportionality of the single emitter decay rate fluctuations with the local density of states fluctuations. Because short-range correlations can alter the single emitter decay rate statistics, light scattering in the near field is crucial for the understanding of correlated random systems.\\
The lack of completeness in the study of emission dynamics extends from simple geometry systems to complex structures. An interesting example is the inclusion of reciprocity symmetry breaking in well known structures, like in emitter-antenna systems. A way of introducing nonreciprocity effects is by applying an external magnetic field and considering the magnetisation properties of the materials that form the system. The effect that involves the modification of the optical properties by the application of a magnetic field is known as magneto-optical effect. 
This is of special importance in hybrid metallic systems that support at the same time plasmon excitations (magnetoplasmonic systems), giving rise to plasmon enhanced magneto-optical activity. 
These structures are particularly interesting because can be applied in biology, sensing and telecommunications. Due to the available fabrication and characterisation technology, it is possible to design and implement complex structures that take into account this magneto-optical effect. Nevertheless, the optimisation processes are directly related to the modelling ability of such systems.\\
The aim of this thesis is to clarify some open questions in magnetoplasmonic systems and random correlated systems. The thesis is organised in two parts.\\
In the first part, we want to understand, through simple physical models, the rich phenomenology obtained in hybrid magnetoplasmonic systems. In particular, we are interested in determine the exact role of interactions between forming elements of the systems and their contribution to the global magneto-optical response.\\
We will also address the antenna effect in the presence of magneto-optical activity. The considered model is an emitter-antenna system in which the antenna supports a plasmon resonance and shows magneto-optical activity.\\
In the second part of this thesis we study the emission lifetime in thin layers of strongly scattering media with non-trivial correlations: search for evidences of peculiar statistical distributions in lifetime statistics (long tail distributions, multimodal, etc). The objective is trying to prove that not only the parameters governing the statistical distributions of lifetimes are non-universal, but also the nature of the distributions themselves depend on the details of the system. \\
Finally, we focus our attention on the light emission dynamics (lifetimes or emission rate statistics) in three-dimensional systems with strong structural correlations undergoing structural phase transitions. The objective is trying to link the statistical properties of emission dynamics with the thermodynamics of the system and prove that there are fundamental differences between lifetime statistics and transport properties.\\





